\chapter{Literature Review}
\label{ch:review}

This chapter establishes the evidence base for the study and evaluates it
critically. It is organised around the questions the field has not settled
rather than around a catalogue of studies, because the design decisions taken
in Chapter~\ref{ch:methodology} are responses to specific unresolved disputes,
not to a consensus.

Four such disputes structure what follows. The first is taxonomic: the field
routinely partitions itself into ``statistical'' and ``machine-learning''
traditions, and Section~\ref{sec:false_dichotomy} argues that this partition
does not carve the problem at its joints. The second concerns the outcome
itself: what counts as sepsis onset, and whether it is a stable enough
object to support the interpretable effect estimates that motivate statistical
modelling (Sections~\ref{sec:label_object} and~\ref{sec:label_action}). The
third concerns validation: what the published external-validation record
actually shows, and why discrimination and clinical utility tell opposite
stories about it (Sections~\ref{sec:generalisation}
and~\ref{sec:auroc_illusion}). The fourth is methodological, and is genuinely
open in the statistical literature rather than merely under-reported: no
survival formulation for ICU event data satisfies all of its own assumptions,
and Section~\ref{sec:survival_dispute} takes a position on which set of
violations is acceptable.

Sections~\ref{sec:comparator_choice} and~\ref{sec:equity_lit} then treat two
narrower questions (which rule-based comparator is worth beating, and
whether subgroup validity has been assessed at all), before
Section~\ref{sec:interpretability_survives} draws the threads together into the
argument that motivates this study's research question, and
Section~\ref{sec:design_warrant} maps each design decision onto the evidence
that warrants it.

A fuller treatment of the same corpus, covering fifty sources including several
whose relevance is contextual rather than design-bearing, is available in the
companion critical review \textit{Statistical and Machine-Learning Approaches
to Real-Time Sepsis Onset Prediction on MIMIC-IV}. The sources that bear on the
decisions taken here are engaged with in this chapter.


\section{Search Strategy}
\label{sec:search_strategy}

The literature search was conducted to identify published work on real-time
sepsis prediction, with particular attention to label sensitivity, validation
methodology, evaluation metrics, and subgroup validity.
Table~\ref{tab:search_strategy} documents the protocol.

\paragraph{What kind of review this is.} This is a \textbf{structured narrative
critical review, not a PRISMA-conformant systematic review}. Searches were
executed to a documented protocol and every included source was verified
against a primary record, but title-and-abstract screening was not conducted
against a full de-duplicated yield, and no second reviewer independently
screened. The distinction matters for the novelty claims in
Section~\ref{sec:contrib}: those claims are stated as ``to the best of our
knowledge'' and are bounded by the coverage this method achieves, which is
deliberate depth on the methodological literature that constrains this study
rather than exhaustive coverage of the MIMIC-IV sepsis corpus. An exhaustive
systematic review could surface a prior study combining the two analyses in
contribution C5; this search did not.

{\small
\begin{longtable}{p{3.2cm}p{9.5cm}}
    \caption{Literature search strategy. This search supports the novelty claims
    for contributions C1 and C5 stated in Section~\ref{sec:contrib}. Because the
    review is narrative rather than systematic, those claims are stated as
    ``to the best of our knowledge'' and are bounded by the coverage above.}\label{tab:search_strategy}\\
    \toprule
    \textbf{Element} & \textbf{Detail} \\
    \midrule
    \endfirsthead
    \toprule
    \textbf{Element} & \textbf{Detail} \\
    \midrule
    \endhead
    \bottomrule
    \endfoot

        Databases searched &
        PubMed/MEDLINE, Embase, Scopus, Web of Science, IEEE Xplore,
        ACM Digital Library, Google Scholar, Cochrane Library,
        ScienceDirect, SpringerLink, Wiley Online Library, Nature
        Portfolio, JAMA Network, BMJ Journals, arXiv/medRxiv \\
        \addlinespace
        Searches executed &
        35 structured searches across the domains listed under ``Search
        terms'' \\
        \addlinespace
        Date range &
        Primary sepsis-prediction literature: January 2015 -- July 2026.
        Methodological and definitional foundations (e.g.\ pooled logistic
        regression, decision-curve analysis) were searched without a date
        limit, and the earliest source retained dates from 1990 \\
        \addlinespace
        Search terms &
        \texttt{(sepsis OR ``sepsis prediction'' OR ``sepsis detection'') AND
            (MIMIC OR ``critical care'' OR ``intensive care'')} combined with each
        of: \texttt{``label definition''}, \texttt{``label sensitivity''},
        \texttt{``onset definition''}, \texttt{``utility score''},
        \texttt{``clinical utility''}, \texttt{``fairness'' OR ``equity'' OR
        ``disparity'' OR ``bias''}, \texttt{``subgroup analysis''},
        \texttt{``competing risks'' OR ``discrete-time survival''} \\
        \addlinespace
        Inclusion criteria &
        (1)~Peer-reviewed journal article, conference paper, or preprint on a
        recognised server (arXiv, medRxiv);
        (2)~Develops or evaluates a real-time (hourly or sub-hourly) sepsis
        prediction model;
        (3)~Uses EHR data from an ICU population;
        (4)~Reports at least one quantitative performance metric \\
        \addlinespace
        Exclusion criteria &
        (1)~Non-English language;
        (2)~Editorials, commentaries, or letters without original results;
        (3)~Models for sepsis severity staging or mortality only (not onset
        prediction);
        (4)~Studies on paediatric or neonatal populations exclusively \\
        \addlinespace
        Verification standard &
        Every included source was checked against a publisher record, PubMed or
        PMC record, or institutional repository entry, with journal, volume,
        year and DOI confirmed. Every numerical claim attributed to a source
        was located in that source's own abstract, results text or
        publisher-hosted record, never inferred from a secondary citation \\
        \addlinespace
        Total included &
        50 sources in the companion critical review (46 newly verified in this
        pass, 4 carried forward from the base review of 19 sources); 45 are
        cited in this thesis \\
        \addlinespace
        Excluded at verification &
        4 items encountered in secondary citation lists could not be verified
        to the standard above and were excluded; 1 further item is retained in
        the companion review with an explicit note of incomplete metadata and
        is not cited in this thesis \\
\end{longtable}
}

Figure~\ref{fig:prisma} summarises the identification, verification and
inclusion process. It follows the reporting logic of PRISMA 2020
\cite{page2021prisma}, an auditable count at every stage, but is
presented as a search-and-verification flow rather than as a PRISMA flow
diagram, because the screening stage a PRISMA diagram reports was not
performed in this form.

\begin{figure}[H]
\centering
\adjustbox{max width=\textwidth}{%
\begin{tikzpicture}[
  node distance=0.55cm,
  box/.style={rectangle, draw=ATUGreen, thick, rounded corners=2pt,
              fill=ATULightGreen!20, text width=7.5cm, align=center,
              font=\small, inner sep=6pt, minimum height=0.8cm},
  excl/.style={rectangle, draw=ATUOrange, thick, rounded corners=2pt,
               fill=ATUSand!40, text width=5.5cm, align=center,
               font=\small, inner sep=5pt, minimum height=0.7cm},
  phase/.style={font=\small\bfseries, text=ATUNavy, rotate=90, anchor=south},
  arr/.style={-{Latex[length=2mm]}, thick, ATUGreen},
  exarr/.style={-{Latex[length=2mm]}, thick, ATUOrange},
]

% --- Identification ---
\node[box] (db) {35 structured searches across\\15 bibliographic sources\\
  (10 methodological domains)};

\node[box, below=of db] (base) {Sources in the base review\\
  ($n = 19$: 18 studies + 1 review)};
\draw[arr] (db) -- (base);

% --- Extension ---
\node[box, below=of base] (cand) {Candidates carried into the\\
  extended critical review};
\draw[arr] (base) -- (cand);

\node[excl, right=1cm of cand] (e2) {Excluded at verification\\
  ($n = 4$): could not be traced\\to a primary record};
\draw[exarr] (cand.east) -- (e2.west);

% --- Verification ---
\node[box, below=of cand] (verif) {Verified against publisher /\\
  PubMed / PMC / repository record\\
  ($n = 46$ newly verified)};
\draw[arr] (cand) -- (verif);

\node[excl, right=1cm of verif] (e3) {
  Retained with caveat ($n = 1$):\\[2pt]
  \begin{minipage}{4.8cm}\raggedright\scriptsize
    incomplete metadata; not cited\\in this thesis
  \end{minipage}
};
\draw[exarr] (verif.east) -- (e3.west);

% --- Included ---
\node[box, below=of verif, fill=ATUGreen!15, draw=ATUGreen!80!black]
  (incl) {Sources in the companion review\\
  ($n = 50$; 4 carried forward)\\
  of which cited in this thesis: 45};
\draw[arr] (verif) -- (incl);

% --- Phase labels ---
\node[phase, left=1.2cm of db]    {Identification};
\node[phase, left=1.2cm of cand]  {Selection};
\node[phase, left=1.2cm of verif] {Verification};
\node[phase, left=1.2cm of incl]  {Included};

\end{tikzpicture}}
\caption{Search, verification and inclusion flow. Counts are auditable at every
  stage, following the reporting logic of PRISMA 2020 \cite{page2021prisma},
  but this is not a PRISMA flow diagram: no title-and-abstract screening pass
  over a full de-duplicated yield was performed, and screening was not
  duplicated by a second reviewer. The review is narrative and critical, and
  the novelty claims it supports are bounded accordingly
  (Section~\ref{sec:contrib}).}
\label{fig:prisma}
\end{figure}

\paragraph{What the search establishes about novelty, and what it does not.}
Two contributions rest on search-based evidence, and both are claims about a
\emph{combination} of analyses rather than about either analysis on its own.
Each component has clear precedent; what the search did not identify is any
study performing them together.

\begin{enumerate}
    \item \textbf{C1 (Label variance under a dynamic survival model).} Label
    sensitivity itself is well established: Johnson et al.\
    \cite{johnson2018comparative}, Cohen et al.\ \cite{cohen2024subtle} and
    Dutta et al.\ \cite{dutta2026performance} all demonstrate it, and the last
    two are the direct precedents for this study. Cohen et al.\ include a
    survival model among the families they compare, so the survival setting is
    not itself new either. What the search did not identify is a study applying
    three or more Sepsis-3 readings to a \emph{discrete-time competing-risks}
    model on MIMIC-IV with the label spread reported as an estimand carrying an
    interval, alongside the model-class spread measured on the same test rows.
    The claim is that combination, not label sensitivity as such.

    \item \textbf{C5 (Performance disparity and label sensitivity across the
    same subgroups).} Subgroup performance disparity in sepsis cohorts is
    reported by Wang, Li, Naidech and Luo \cite{wang2022equity}, and subgroup
    variation in \emph{who gets labelled} is reported in the same paper. It is
    important to be exact about what that study is: it predicts
    \textbf{mortality among patients already identified as septic}, not sepsis
    onset. It therefore cannot serve as the prior art against which an onset
    model's subgroup performance is judged, and the search did not identify any
    onset-prediction study reporting subgroup discrimination at all. Studies
    examining label sensitivity for onset prediction (Cohen et al., Dutta et
    al.) do not stratify by demographic subgroup. The novel element is the
    conjunction (subgroup discrimination \emph{and} subgroup label
    sensitivity, measured on the same subgroups, in the same cohort, under
    common false-discovery control), and not either half.
\end{enumerate}

\paragraph{A bounded claim about clinical utility.}
Wang et al.\ \cite{wang2025srpm} reviewed ninety-one real-time sepsis
prediction models. Of these, only a minority reported the PhysioNet Utility
Score or an equivalent net-benefit metric; the majority reported only AUROC
and/or AUPRC. The claim that published models do not demonstrate good clinical
utility is therefore based on the subset for which a utility or net-benefit
metric was available: among those, the median external Utility Score was
$-0.164$, indicating net harm. For models reporting only AUROC, clinical
utility could not be assessed, and no claim about their utility is made here.

\subsection*{Key Concepts}

Several domain-specific terms recur throughout this chapter:

\begin{itemize}
    \item \textbf{MIMIC-IV} (Medical Information Mart for Intensive Care, version
    IV): a freely available, de-identified database of approximately 74,000 ICU
    admissions at Beth Israel Deaconess Medical Center, Boston, covering
    2008--2019 \cite{johnson2023mimiciv}.
    \item \textbf{Sepsis-3}: the current international consensus definition of
    sepsis as life-threatening organ dysfunction caused by a dysregulated host
    response to infection, operationalised through a combination of suspected
    infection (antibiotic order plus culture draw) and a rise in the Sequential
    Organ Failure Assessment (SOFA) score of $\geq$2 points
    \cite{singer2016sepsis3}.
    \item \textbf{AUROC} (Area Under the Receiver Operating Characteristic
    curve): a discrimination metric that measures a model's ability to rank
    positive cases above negative cases, ranging from 0.5 (random) to 1.0
    (perfect separation). It is invariant to any monotone transformation of the
    score, and therefore cannot see miscalibration; it does not account for
    alert timing or false-alarm cost.
    \item \textbf{Utility Score}: a net-benefit metric proposed by the
    PhysioNet/Computing in Cardiology Challenge \cite{reyna2020utility} that
    rewards early, true alerts and penalises late alerts and false alarms,
    producing values that can be negative when a model causes more harm than
    benefit.
    \item \textbf{NEWS2} (National Early Warning Score 2): a standardised
    clinical scoring system aggregating six physiological parameters (respiratory
    rate, oxygen saturation, systolic blood pressure, pulse rate, level of
    consciousness, and temperature) into a single risk score used for bedside
    deterioration screening.
\end{itemize}

\subsection*{Research Questions for the Review}

The review is structured around four questions that directly inform the
study's design:

\begin{enumerate}
    \item[\textbf{RQ1.}] How much does the choice of sepsis label definition
    affect model performance, and is this effect comparable in magnitude to
    the choice of model architecture?
    \item[\textbf{RQ2.}] What methodological pitfalls (validation strategy,
    metric selection, treatment-time leakage)most distort reported
    performance in the sepsis prediction literature?
    \item[\textbf{RQ3.}] What statistical framework is appropriate for
    modelling hourly ICU event data when the standard proportional-hazards
    assumption is violated?
    \item[\textbf{RQ4.}] Does the existing equity literature on sepsis
    prediction account for label-level variation across demographic subgroups,
    or only model-level variation?
\end{enumerate}


\section{The Field's Usual Partition Is Not the Load-Bearing One}
\label{sec:false_dichotomy}

Reviews of this literature conventionally divide it into a classical
statistical tradition and a machine-learning tradition, and the division
usually carries an implicit evaluative charge: interpretable models on one
side, accurate models on the other. The evidence does not support the
partition, and treating it as real leads to a specific analytical error that
this study is designed to avoid.

The boundary cases dissolve it. The Artificial Intelligence Sepsis Expert
\cite{nemati2018aise}, routinely classified as a machine-learning system, is a
modified Weibull--Cox proportional-hazards formulation: the leading
``machine-learning'' sepsis model of its generation is built on a
survival-analytic backbone. TREWScore \cite{henry2015trewscore}, a Cox model
and the canonical \emph{statistical} benchmark for this task, is described
throughout the machine-learning literature as an early machine-learning
early-warning system, and its commercial successor is marketed as one
\cite{adams2022trews}. Cohen et al.\ \cite{cohen2024subtle} place tree-based,
deep-learning and survival models in a single comparison frame because the
distinction is not load-bearing for the question they ask. And TRIPOD+AI
\cite{collins2024tripodai}, the governing reporting standard, covers
``regression or machine learning'' under one checklist, having concluded that
the reporting failures are common to both.

The partition that \emph{is} load-bearing separates models by what they assume
about time.

\begin{itemize}
    \item \textbf{Static models} take a snapshot (baseline, or day one, or
    the worst value in the first 24 hours), and issue one prediction. The
    great majority of the MIMIC-IV statistical corpus sits here: a candidate
    biomarker or composite index is regressed on mortality through an almost
    invariant pipeline of variable selection, Cox and logistic regression,
    splines, a nomogram, and ROC, calibration and decision curves. These are
    \emph{prognostic} instruments and cannot in principle perform real-time
    onset detection.
    \item \textbf{Dynamic models} consume a stream and update a risk estimate as
    it arrives. TREWScore, the Artificial Intelligence Sepsis Expert, the
    PhysioNet Challenge entries \cite{reyna2020utility}, Moor et al.'s deep
    temporal model \cite{moor2023review}, and landmark and joint models all sit
    here. Whether the update rule is a partial likelihood, a recurrent network
    or a boosted hazard is a secondary question.
\end{itemize}

Reframing the taxonomy this way corrects a claim that the statistics-versus-ML
framing tends to generate: that statistical methods have neglected onset
prediction. They have not. Henry et al.\ \cite{henry2015trewscore} fitted a Cox
model for septic-shock onset on MIMIC-II; Fagerstr{\"o}m et al.\
\cite{fagerstrom2019lisep} fitted one on MIMIC-III using Henry et al.'s exact
inputs and target as the baseline for an LSTM; and Cohen et al.'s survival arm
does so on MIMIC-III again. What is genuinely absent is narrower and harder: a
properly specified, assumption-tested, label-sensitivity-analysed, temporally
validated dynamic model for \emph{onset} on MIMIC-IV.

The Fagerstr{\"o}m result also bounds a common claim about model class.
Christodoulou et al.\ \cite{christodoulou2019cox}, reviewing 71 studies drawn
from 927 screened, found no performance advantage for machine learning over
logistic regression in clinical prediction, and observed potential bias in the
validation procedures of 68\% of them. That is the strongest available warrant
for treating a classical specification as a serious contender. But its scope is
binary outcomes on tabular data, and sepsis onset prediction is not that
problem: it is a temporal problem over irregularly sampled multivariate
streams in which the quantity of interest is a lead time. On precisely that
axis Fagerstr{\"o}m et al.\ report the deep model detecting patients up to
twenty hours earlier than the Cox model at comparable sensitivity and
specificity. The defensible position is therefore narrower than ``statistics
matches machine learning'': it is that the equivalence is established for
static risk stratification and is contested for the temporal early-warning
task. H6 in this study is a direct test of that narrower proposition, on one
dataset, under one label, with the split and the evaluation held fixed: the
conditions Christodoulou et al.\ find are usually absent.

Two further consequences follow for design. Moor et al.'s systematic review
\cite{moor2021review} identifies low comparability and reproducibility as the
field's principal defect, and singles out case--control matching as a concrete
trap: when a control patient's ``onset time'' is chosen near discharge while a
case is aligned near onset, the two are trivially separable and apparent
performance is inflated. Fleuren et al.\ \cite{fleuren2020ml} reach a
compatible conclusion from a meta-analysis of diagnostic test accuracy,
attributing the between-study heterogeneity that prevents pooling to
methodological rather than clinical variation: heterogeneity that Cohen et
al.\ later show is partly label variation. This study therefore performs no
case--control matching at all: every at-risk person-hour contributes, which
removes the alignment decision from the analyst's hands.


\section{The Sepsis Label Is Not a Single Object}
\label{sec:label_object}

The common assumption that ``the MIMIC-IV sepsis cohort'' is one well-defined
group is wrong, and the evidence against it is now strong enough to treat as an
established property of the problem rather than a curiosity.

Johnson et al.\ \cite{johnson2018comparative} applied five retrospective sepsis
identification algorithms, including the Sepsis-3 clinical criteria, to
11,791 qualifying ICU admissions drawn from 23,620, and found agreement across
them that the authors themselves characterise as only acceptable (Cronbach's
$\alpha$ = 0.40--0.62). These algorithms are the very queries that MIMIC-based
studies run to construct their cohorts. Different studies claiming to study the
same disease in the same database are therefore, measurably, not studying the
same patients.

Cohen et al.\ \cite{cohen2024subtle} established what this costs a prediction
model, and their design is the direct precedent for this thesis. Working in
MIMIC-III, they held the data fixed and varied only the reading of the
Sepsis-3 onset rule across three defensible interpretations of the same
consensus definition. The number of identified septic admissions ranged from
867 to 2,178, a factor of 2.5, purely as a function of that reading. They
then trained representative models spanning tree-based, deep-learning and
survival families under each definition. Holding the definition fixed, the best
method gained 1--5\% AUROC over its rivals; holding the method fixed and
varying the onset definition produced swings of 0--6\%. The analyst's
discretionary reading of the label moved performance at least as much as the
entire choice of modelling paradigm.

Dutta et al.\ \cite{dutta2026performance} confirm this is not a MIMIC artefact,
and the corroboration is unusually strong because almost nothing about the
setting is shared. A locally trained gradient-boosted model generating
predictions every fifteen minutes across nine acute-care hospitals was
evaluated on 198,494 encounters against three electronically computable
definitions: Sepsis-3, the SEP-1 management bundle, and the CDC Adult Sepsis
Event. Performance varied substantially by definition, with expected
calibration error differing across them (0.05 for Sepsis-3, 0.06 for Adult
Sepsis Event, 0.07 for SEP-1) and decision-curve analysis showing higher
standardised net benefit under Sepsis-3 and SEP-1. Cohen et al.\ established
definitional sensitivity retrospectively, in MIMIC-III, at $n \approx 2{,}000$,
with research models; Dutta et al.\ reproduce it prospectively, in a live
deployment, at $n \approx 200{,}000$, with a commercial model. The finding is
robust to dataset, scale, model class and setting.

The critical gap in this line of work is what makes it the starting point for
this study rather than its conclusion. Both papers demonstrate \emph{that} the
label matters; neither separates the part of the label effect that is
guaranteed by the definitions from the part discovered in the data, and neither
attaches an interval to the comparison between label variance and model
variance. A spread of point estimates is not an estimand. Cohen et al.'s
``0--6\% versus 1--5\%'' is a comparison of two ranges, each computed without
uncertainty, on a cohort of roughly two thousand admissions. Whether the two
sources of variance are genuinely of comparable size is therefore not settled
by it, and Sections~\ref{sec:label_anchor_results} and~\ref{sec:decomposition}
of this thesis are constructed to settle it on a cohort thirty times larger
with intervals attached to both quantities.

Figure~\ref{fig:label_variance_concept} illustrates the mechanism: the same
patient population, subjected to different but equally valid Sepsis-3
interpretations, produces overlapping but distinct cohorts.

\begin{figure}[H]
    \centering
    \begin{tikzpicture}[
        every node/.style={font=\small},
        cohort/.style={draw, circle, minimum size=2.8cm, thick, fill opacity=0.15},
    ]
    \node[cohort, fill=ATUGreen, draw=ATUGreen] (a) at (0, 0) {};
    \node[cohort, fill=ATUOrange, draw=ATUOrange] (b) at (1.6, 0) {};
    \node[cohort, fill=ATUNavy, draw=ATUNavy] (c) at (0.8, 1.3) {};

    \node[below=0.2cm, font=\footnotesize\bfseries, text=ATUGreen] at (a.south) {Reading 1};
    \node[below=0.2cm, font=\footnotesize\bfseries, text=ATUOrange] at (b.south) {Reading 2};
    \node[above=0.2cm, font=\footnotesize\bfseries, text=ATUNavy] at (c.north) {Reading 3};

    \node[font=\scriptsize] at (0.8, 0.4) {\textbf{Shared}};

    \draw[decorate, decoration={brace, amplitude=5pt, raise=3pt}]
    (-1.6, -2.0) -- (3.2, -2.0)
    node[midway, below=10pt, font=\footnotesize, align=center]
        {Same ICU population, same EHR data\\different ``sepsis cohorts''};

    \node[draw, rounded corners, fill=white, inner sep=6pt,
        font=\footnotesize, align=left] at (6.9, 0.5) {
        \textbf{Reported by Cohen et al.}\\
        \textbf{(MIMIC-III, three readings):}\\[3pt]
        Cohort size: 867--2{,}178\\
        \quad admissions\\[3pt]
        Best method vs.\ rivals:\\
        \quad $+1$--$5\%$ AUROC\\[2pt]
        Onset definition varied:\\
        \quad $0$--$6\%$ AUROC
    };
    \end{tikzpicture}
    \caption{Schematic of label variance. Three valid Sepsis-3 interpretations
    applied to the same ICU population produce overlapping but distinct
    cohorts. The boxed quantities are those \emph{reported by Cohen et al.}\
    \cite{cohen2024subtle} on MIMIC-III; they are shown to convey the scale of
    the effect that motivates this study and are not predictions, targets, or
    results of the present work. The corresponding measurements on MIMIC-IV are
    reported in Chapter~\ref{ch:evaluation}.}
    \label{fig:label_variance_concept}
\end{figure}


\section{The Label Is Constituted by Clinical Action}
\label{sec:label_action}

Sepsis-3 defines sepsis as organ dysfunction in a patient with suspected
infection \cite{singer2016sepsis3}, and the operational criteria derived by
Seymour et al.\ \cite{seymour2016assessment} render ``suspected infection''
through two clinician actions: the drawing of a culture and the administration
of an antibiotic. The consequence is a labelling rule that reads two clocks,
one of which is a record of what the treating team did.

Kamran et al.\ \cite{kamran2024evaluation} pursued this to its conclusion and
produced the deepest threat to the enterprise. Clinicians frequently recognise
and begin treating sepsis before the formal criteria are met. A model trained
to predict the moment the criteria are met is therefore partly trained to
predict \emph{when the clinician will act}, and the physiological record
available at prediction time is itself shaped by that developing suspicion:
tests are ordered, monitoring escalates, and observations are recorded that
would not otherwise exist. Re-evaluating the Epic Sepsis Model on 77,582
hospitalisations, the authors report an AUROC of 0.62 when predictions issued
after clinical recognition are retained, falling to 0.47, worse than chance, once they are excluded.

For a hazard model the implication is exact. A coefficient estimated on the
interval up to Sepsis-3 onset is in part a coefficient for clinician behaviour.
Lactate is drawn because someone is worried; the predictor and the outcome
share a common cause in the clinician's suspicion, and no amount of
proportional-hazards testing detects this, because the model is correctly
estimating the hazard of a contaminated event. The remedy is definitional
rather than statistical: evaluation must be anchored to a prediction time that
precedes clinical action.

Lauritsen et al.\ \cite{lauritsen2021framing} generalise the problem beyond the
label to the surrounding design. They formalise \emph{framing} (the
observation window, the prediction window, the window shift, and the event that
triggers a prediction) enumerate eight possible framing structures, and
evaluate four on the same underlying data. Framing determines both
discrimination and calibration, and, critically, can produce opposing
interpretations of the same physiological variable: their worked example is
SpO$_2$, whose apparent relationship to sepsis risk reverses direction with the
framing. A model with strong discrimination \emph{and} strong calibration may
nonetheless be clinically unusable.

Taken together these two results close a trap around the argument for
interpretable models. If framing can flip the sign of a coefficient, then the
coefficient reports a property of the study design as much as a property of the
patient. An uninterpretable model that is wrong is a nuisance; an
\emph{interpretable} model that is confidently wrong about the direction of a
physiological effect is a hazard, because clinicians act on directions. H2 in
this study is the direct test of whether that hazard is realised on MIMIC-IV.

Where this study departs from Kamran et al.\ is in the scope of the remedy.
Treatment-anchored evaluation restricts \emph{when} performance is measured,
but it leaves the anchoring inside the label untouched: a Sepsis-3 case is
still a case only because someone ordered a culture and an antibiotic. A
comparison confined to Sepsis-3 readings (or to SEP-1 and the CDC Adult
Sepsis Event, both of which are defined on antibiotic days) cannot separate
definitional from empirical label variance, because every member of that family
keys on treatment timing. Breaking the circularity requires a label built from
physiology alone, which is why this study constructs one
(Section~\ref{sec:alt_labels}) rather than confining itself to the consensus
family.


\section{The Generalisation Gradient}
\label{sec:generalisation}

Assembling the validation evidence into one frame exposes a consistent
structure that no single study makes visible.
Table~\ref{tab:generalisation} collects it.

\begin{table}[H]
    \centering\small
    \begin{tabularx}{\textwidth}{p{2.6cm}p{2.5cm}Xp{2.9cm}}
    \toprule
    \textbf{Study} & \textbf{Model} & \textbf{Transition} &
    \textbf{Degradation} \\
    \midrule
    Nemati et al.\ \cite{nemati2018aise} & Weibull--Cox (AISE) &
    Emory $\rightarrow$ MIMIC-III; AUROC 0.83--0.85 &
    Reported indistinguishable \\
    \addlinespace
    Reyna et al.\ \cite{reyna2020utility} & 853 Challenge entries &
    Two seen hospital systems $\rightarrow$ one unseen; rank correlation
    $\rho = 0.949$ between seen systems &
    $\rho = -0.033$ and $0.013$ to the unseen system \\
    \addlinespace
    Moor et al.\ \cite{moor2023review} & Deep temporal &
    Four international ICU databases, 136,478 admissions; internal AUC 0.846 &
    External AUC \textbf{0.761} ($-0.085$); 0.807 with 10\% fine-tuning \\
    \addlinespace
    Wang et al.\ \cite{wang2025srpm} & 91 models (meta-level) &
    Internal $\rightarrow$ external, full-window &
    AUROC $0.811 \rightarrow 0.783$ (n.s.); \textbf{Utility
    $0.381 \rightarrow -0.164$} \\
    \addlinespace
    Wong et al.\ \cite{wong2021epic} & Epic Sepsis Model v1 &
    Proprietary $\rightarrow$ Michigan Medicine, 38,455 hospitalisations &
    AUROC \textbf{0.63}; PPV 12\% \\
    \addlinespace
    Ostermayer et al.\ \cite{ostermayer2024external} & Epic Sepsis Model v1 &
    Proprietary $\rightarrow$ two county emergency departments &
    Sensitivity \textbf{14.7\%}; PPV 7.6\% \\
    \addlinespace
    Wong et al.\ \cite{wong2026epic} & Epic Sepsis Model v2 &
    Prospective, four US health systems, 227,091 encounters &
    AUROC 0.82--0.92 but \textbf{PPV 0.13--0.26}; 21--35 alerts per case \\
    \addlinespace
    Guo et al.\ \cite{guo2022temporal} & Feed-forward network &
    MIMIC-IV, within database, across admission eras &
    Sepsis \textbf{worst-degraded} of four tasks \\
    \bottomrule
    \end{tabularx}
    \caption{The published external-validation record for real-time sepsis
    prediction. Nemati et al.'s internal--external equivalence is the exception
    and is instructive: their external set shares an era, a country and a care
    model with the development set, which is a mild shift by the standards of
    Moor et al.'s Netherlands-and-Switzerland comparison. The apparent
    counter-example is better read as evidence that the size of the external
    penalty scales with the distance of the shift.}
    \label{tab:generalisation}
\end{table}

Three regularities follow, and each dictates a design decision taken in
Chapter~\ref{ch:methodology}.

\paragraph{The external penalty for a well-built model is roughly 0.085 AUROC.}
This is Moor et al.'s 0.846 $\rightarrow$ 0.761 across four databases spanning
the United States, the Netherlands and Switzerland, and it is consistent in
direction, if not in significance, with Wang et al.'s meta-level 0.811
$\rightarrow$ 0.783. The number is specific enough to pre-specify, which is
what H5 does; and a study observing a substantially \emph{smaller} drop should
first suspect that its external cohort is insufficiently distant rather than
conclude that its model is unusually robust.

\paragraph{Rank order does not transfer.} Reyna et al.\ report that team
Utility Scores were strongly rank-correlated between the two hospital systems
participants had trained on ($\rho = 0.949$) but essentially uncorrelated with
performance on the third, unseen system ($\rho = -0.033$ and $0.013$). The
ranking of algorithms carried almost no information out of distribution. For a
study that intends to compare its model against published competitors this is
foundational: a benchmark is informative only when run on the same data, under
the same label, with the same split, and even then it licenses no inference
about which model would win at a different hospital. Every comparator in this
study is therefore refitted and rescored inside the same pipeline, and no
performance figure is compared against a published one except to bound an
expectation.

\paragraph{The split design is not neutral, and sepsis is the worst case.}
Guo et al.\ \cite{guo2022temporal} partitioned MIMIC-IV into admission eras
(2008--2010, 2011--2013, 2014--2016, 2017--2019) and evaluated four prediction
tasks under temporal shift: mortality, long length of stay, sepsis, and
invasive ventilation. Sepsis was the worst-affected. Domain-generalisation and
unsupervised domain-adaptation algorithms failed to beat plain empirical risk
minimisation (AUROC differences $-0.003$ to $0.050$), so there is at present no
algorithmic fix. The consequence is unavoidable and task-specific: a random
train/test split on MIMIC-IV is optimistically biased for sepsis, and more so
than for the neighbouring ICU outcomes a reader might use to calibrate their
scepticism. This study's primary validation is therefore a temporal split; the
random split is retained only to quantify the optimism, which is H4.

A qualification the field does not usually state: Dutta et al.\
\cite{dutta2026performance} did \emph{not} detect temporal drift over their
eighteen-month prospective period. Shift over eighteen months inside one health
system is a different phenomenon from shift across a decade of MIMIC-IV, and
the two results are compatible. What follows is that ``temporal validation''
is not a single practice, and the era gap matters as much as the fact of
splitting by time.


\section{Discrimination and Clinical Value Tell Opposite Stories}
\label{sec:auroc_illusion}

The most consequential comparison in this literature sits inside a single
paper. Wang et al.\ \cite{wang2025srpm} report, across ninety-one models, that
the median AUROC moved from 0.811 internally to 0.783 externally, a
difference they report as not statistically significant, while the median
Utility Score moved from $+0.381$ to $-0.164$. The same models, the same
transition, and two metrics that disagree about whether anything happened. On
the Challenge's utility calculus, a negative score means that deploying the
median externally validated sepsis model is worse than deploying nothing.

The mechanism is not mysterious. AUROC is a ranking statistic, invariant to
monotone transformation of the score, so it cannot see miscalibration; and
under low prevalence it is dominated by the enormous population of easily
classified true negatives. What changes across sites is not primarily the
\emph{ordering} of patients but the \emph{location} of the score distribution
relative to the alerting threshold, which is exactly what determines the
false-alarm rate, and therefore the utility. Van Calster et al.\
\cite{vancalster2019calibration} name calibration the Achilles heel of
predictive analytics for this reason: any model driving a threshold-based
action is a calibration-dependent instrument whatever its AUROC.

Two deployment results make the same point visible in a second way. Wong et
al.\ \cite{wong2026epic}, validating the Epic Sepsis Model v2 prospectively
across 227,091 encounters at four health systems, report encounter-level AUROC
between 0.82 and 0.92, an enormous improvement on version 1's 0.63
\cite{wong2021epic}, alongside positive predictive value of only 0.13 to
0.26, requiring clinicians to investigate 21 to 35 alerts to find one case.
Discrimination near ceiling; usability still contested. Ostermayer et al.\
\cite{ostermayer2024external}, validating version 1 independently across two
county emergency departments, report sensitivity of 14.7\% and PPV of 7.6\%,
worse than Wong et al.'s original finding, in a different setting and case mix.

The v2 result deserves a qualification, because it cuts against the sceptical
reading. Version 1 was a logistic regression and version 2 is a gradient-boosted
ensemble, so the improvement is evidence against the strongest form of
Christodoulou et al.'s claim \cite{christodoulou2019cox}. But version 2 also
changed from globally trained to locally retrainable, and Moor et al.\ have
already shown that local fine-tuning recovers roughly half of the external
penalty ($0.761 \rightarrow 0.807$). Model class and localisation are therefore
confounded in the published comparison, and the evidence does not permit them to
be separated. What is not ambiguous is that PPV remained low and alert burden
high at AUROC 0.92.

Two benchmarks carry forward from this section into the evaluation plan. Moor
et al.\ report their deployed system raising \textbf{1.4 false alerts per true
alert} while detecting 80\% of septic patients a median of 3.7 hours (95\% CI
3.0--4.3) before onset. That is the reference burden against which any alert
volume should be read, and the 3.7-hour figure is the honest lead time of a
good modern system, not the 28 hours of the original TREWScore headline
\cite{henry2015trewscore}. And Seymour et al.\ \cite{seymour2017time}, across
49,331 patients at 149 New York hospitals, estimate an odds ratio of 1.04 per
additional hour to antibiotic administration. A model delivering Moor et al.'s
3.7-hour lead time with perfect clinician response would therefore purchase an
improvement of order $1.04^{3.7} \approx 1.16$ on the odds scale: real but
modest, available only if the alert is acted on, only if the patient would not
otherwise have been recognised, and only after the false alerts have been paid
for.

The consequence for this study's evaluation plan is direct. A study reporting
AUROC, sensitivity, specificity and lead time, and concluding that its model
generalises, would have failed to measure the thing that fails. Calibration,
AUPRC, the Utility Score, alert burden and decision-curve analysis are not
enhancements to that list; they are the list. This is why the Utility Score is
the primary metric here and AUROC is descriptive, and, given that the
Challenge metric is threshold-dependent, why it is reported across a swept grid
of operating points rather than at a single cut-off
(Section~\ref{sec:clinical_utility}).


\section{The Survival-Methods Dispute Has No Clean Resolution}
\label{sec:survival_dispute}

A study proposing a survival model for ICU event data inherits a dispute that
the applied literature usually passes over. It is worth stating in full,
because the specification adopted in Chapter~\ref{ch:methodology} is a position
in it rather than a default.

\begin{enumerate}
    \item \textbf{Proportional hazards frequently fails in sepsis cohorts.}
    Kasal et al.\ \cite{kasal2004cox} compared a standard Cox model against
    Gray's flexible survival model in a large severe-sepsis cohort and found
    hazards for death after sepsis frequently non-proportional. This is an
    under-cited caution: the assumption is not merely untested in most MIMIC
    work, it is known to fail in this disease.

    \item \textbf{Independent censoring also fails.} Resche-Rigon et al.\
    \cite{rescherigon2006competing} argue that discharge alive is not a
    non-informative censoring event in ICU cohorts: patients discharged alive
    are systematically unlike those who remain, so the assumption underlying
    Kaplan--Meier estimation is violated. The standard remedy is the Fine--Gray
    subdistribution hazard model \cite{finegray1999subdistribution}.

    \item \textbf{The remedy is itself contested.} Schoenfeld
    \cite{schoenfeld2006survival} argues that survival methods, explicitly
    including competing-risks methods, are inappropriate for ICU outcome
    studies, and that studies seeking prognostic covariates should use logistic
    regression instead.

    \item \textbf{Landmarking does not restore proportional hazards.} The
    landmark approach, refitting at each prediction time on the subset still
    at risk, is the pragmatic route to a continuously updated score
    \cite{vanhouwelingen2007landmark}, and is what most proposals in this space
    gesture toward. It does not in general satisfy the proportional-hazards
    assumption. Landmarking is a different model with its own approximation
    error, not an escape.

    \item \textbf{Joint models restore rigour at a prohibitive cost.} Modelling
    the longitudinal biomarker trajectories and the event process through shared
    latent structure outperforms landmarking when the longitudinal sub-model is
    correctly specified \cite{rizopoulos2017joint}, but the approach is
    computationally expensive, requires a common baseline time, and scales
    poorly to the heterogeneous, irregularly sampled, high-dimensional streams
    characteristic of an ICU record.
\end{enumerate}

There is no assumption-free route. What there is, and what this study adopts,
is a formulation that dissolves several of these problems rather than testing
them: the \textbf{discrete-time hazard model}, fitted as a pooled logistic
regression over person-hour records with time-varying covariates and a flexible
baseline hazard. D'Agostino et al.\ \cite{dagostino1990pooled} established that
pooled logistic regression is mathematically equivalent to time-dependent Cox
regression when the intervals are short and per-interval event rates are low,
and Ngwa et al.\ \cite{ngwa2016comparison} confirmed the equivalence by
simulation. One-hour ICU intervals with hourly sepsis probability below 1\%
satisfy both conditions comfortably.

Its properties answer the dispute point by point. It is native to the data's
structure, because ICU data is already binned hourly by convention
\cite{reyna2020utility, nemati2018aise}: a discrete-time hazard model
\emph{is} an hourly risk score, and no reconciliation between the model's time
scale and the deployment's is required. It does not require proportional
hazards, so item (1) is dissolved rather than tested. It accommodates competing
risks directly through a multinomial cause-specific formulation over the
person-hour \cite{rescherigon2006competing, heyard2019dynamic}, so that
discharge alive and death without sepsis are modelled as what they are, which
answers item (2) without incurring item (3). It answers Schoenfeld's objection
in item (3) on its own terms, being a logistic regression, while retaining the
time-to-event structure he would discard. Its coefficients are log-odds of
onset in the next hour given survival sepsis-free to this hour: arguably a
more clinically legible quantity than a hazard ratio. And it emits exactly the
object the Challenge utility metric consumes, an hourly risk score, so it is
directly comparable with the machine-learning literature. Heyard et al.\
\cite{heyard2019dynamic, heyard2020evaluation} supply the corresponding
development and validation methodology for discrete time-to-event data with
competing risks.

The Cox model is retained in this study, but as a deliberately weak comparator
rather than as the primary specification: its role is to quantify what the
field's default choice costs, and its Schoenfeld test is reported as a finding
about the default rather than as a check on the model actually used.


\section{Which Rule-Based Comparator Is Worth Beating}
\label{sec:comparator_choice}

Most MIMIC-IV sepsis models benchmark against SIRS, qSOFA and NEWS. Two results
constrain how that benchmark should be read.

The Surviving Sepsis Campaign guidelines \cite{evans2021ssc} issue a strong
recommendation \emph{against} using qSOFA as a sole screening tool, on the
grounds of poor sensitivity. Demonstrating superiority over a score the
governing clinical guideline has already told practitioners not to use for
screening is a weak claim, and a study whose headline is that it beats qSOFA
has chosen an opponent that has already conceded. NEWS2 is therefore the minimum
meaningful rule-based comparator, and this study reports SIRS and qSOFA for
completeness while treating NEWS2 as the benchmark.

Mellhammar et al.\ \cite{mellhammar2020news2} supply the second constraint, and
it is a direct precedent for the enterprise rather than a background caution.
They derived a novel score using LASSO on prospectively collected emergency
department data and validated it against NEWS2 and RETTS; the derived score did
not outperform NEWS2 (AUC 0.73 versus 0.69, $p = 0.32$), and the authors
conclude that even with a statistical approach they could not construct better
risk stratification scores for sepsis than NEWS2. Deriving a statistical score
and benchmarking it against rule-based scores has been tried, and it did not
work. This is a valuable negative result and it shapes what this study claims:
the contribution here is not a better score, and the comparison against NEWS2 is
used to locate the model within a family of estimators rather than to argue for
deployment.


\section{Equity and Subgroup Validity}
\label{sec:equity_lit}

Wang, Li, Naidech and Luo \cite{wang2022equity} provide the only source located
by this search that examines subgroup validity in a MIMIC sepsis cohort, and it
is important to characterise it precisely, because its scope is narrower than
the citation is often taken to imply.

The study \cite{wang2022equity} analysed 11,791 MIMIC-III critical-care
patients, of whom 5,783 met Sepsis-3 criteria, and reports two findings. First, six competing sepsis
identification criteria selected systematically different sub-populations with
respect to race, marital status, insurance type and language: who counts as
septic in MIMIC is not demographically neutral. Second, \textbf{mortality}
prediction performance degraded significantly when a model trained on the full
population was applied to Asian patients, Hispanic patients and
Spanish-speaking patients, with pairwise testing detecting significant
discrepancies between Asian and White patients and between English- and
Spanish-speaking patients.

The second finding concerns \textbf{mortality prediction among patients already
identified as septic}, not sepsis onset prediction. It is therefore not a
result about the task studied here, and it cannot serve as the prior estimate
against which an onset model's subgroup discrimination is compared. What it
does establish is the conjunction that motivates this study's equity analysis:
the label is differentially sensitive by subgroup, \emph{and} model performance
is differentially poor for overlapping subgroups. Subgroup analysis is
therefore not a robustness check to be appended in a sensitivity section; it is
a validity requirement, and it is entangled with the label problem of
Section~\ref{sec:label_object}.

The gap is that the two halves have never been measured together on the same
subgroups in the same cohort. Wang et al.\ report both, but for different
targets (label composition for sepsis identification, discrimination for
mortality), so the two cannot be placed side by side. Studies examining label
sensitivity for onset prediction do not stratify demographically at all. This
study reports, for each subgroup, both the variation in who is labelled septic
under each label variant and the variation in onset-model discrimination, under
common false-discovery control, so that the relative precision of the two
analyses is visible rather than assumed.

One methodological caution follows from the same source and governs how the
results in Section~\ref{sec:equity} must be read. MIMIC's demographic fields
are administrative records, not measured population attributes. The race field
mixes broad categories with more granular ones drawn from the same population,
and includes values that record the absence of an answer rather than an answer.
Subgroup contrasts computed on such a field compare administrative strata, and
the interpretation has to be stated in those terms.


\section{Does Interpretability Survive the Label Problem?}
\label{sec:interpretability_survives}

The case for a statistical specification in this domain is epistemic rather
than performance-based: coefficients are interpretable, assumptions are
checkable, and the model can be transparently benchmarked. Christodoulou et
al.\ \cite{christodoulou2019cox} supply the supporting premise that little
discriminative performance is sacrificed by the choice.

Cohen et al.\ \cite{cohen2024subtle} and Lauritsen et al.\
\cite{lauritsen2021framing} attack that case at its root. If the identified
cohort varies by a factor of 2.5 with the reading of the onset rule, and if
problem framing can reverse the sign of a covariate's estimated association,
then a coefficient reported without an accompanying label-sensitivity analysis
is not an interpretable quantity. It is a number whose provenance is partly
biological and partly bureaucratic, and the reader is given no means of
decomposing it.

This does not defeat the case for interpretable models. It relocates it. The
interpretability of a statistical model is a \emph{conditional} property: it
holds given a stable and explicitly specified estimand, and the field has not
supplied one. The discipline required to supply one is precisely the discipline
a classical specification is better equipped to impose. A coefficient shown to
be stable across three defensible readings of the Sepsis-3 onset rule is a
genuinely interpretable quantity; one shown to \emph{flip} across those
readings is a genuinely important negative result. Either way the analysis is
informative, and it is an analysis this search did not find performed for a
survival model on MIMIC-IV.

That observation is the seed of the research question in
Section~\ref{sec:rq}. The field has established that the label matters
(Section~\ref{sec:label_object}), that it is constituted by clinical action
(Section~\ref{sec:label_action}), that the split design is not neutral
(Section~\ref{sec:generalisation}), and that the metric conventionally reported
cannot see the failure (Section~\ref{sec:auroc_illusion}). Each of these has
been demonstrated separately, in different databases, under different models,
by different groups. What has not been done is to hold one database and one
pipeline fixed and ask how much of the reported performance each of these
analyst decisions accounts for, relative to the choice of model that the
literature spends most of its effort on. That question is answerable, it is
answerable with the intervals the constituent studies lack, and it is what this
thesis does.


\section{Summary: Design Decisions and Their Evidence}
\label{sec:design_warrant}

Table~\ref{tab:design_warrant} maps each major design decision onto the
evidence that warrants it and the section in which that evidence is evaluated.

{\small
\begin{longtable}{p{3.1cm}p{5.6cm}p{4cm}}
    \caption{Design decisions and their evidence warrants.}\label{tab:design_warrant}\\
    \toprule
    \textbf{Decision} & \textbf{Reason} & \textbf{Source} \\
    \midrule
    \endfirsthead
    \toprule
    \textbf{Decision} & \textbf{Reason} & \textbf{Source} \\
    \midrule
    \endhead
    \bottomrule
    \endfoot

        Three label variants & Label variance is of the same order as model
        variance, on the only comparison available &
        Cohen et al.\ \cite{cohen2024subtle} \\
        Label as variable & Variance is a finding, not a nuisance; confirmed
        prospectively at scale &
        Cohen; Dutta et al.\ \cite{dutta2026performance} \\
        Discrete-time model & Cox-equivalent at short intervals; no PH
        assumption; answers Schoenfeld's objection &
        D'Agostino \cite{dagostino1990pooled};
        Ngwa \cite{ngwa2016comparison};
        Schoenfeld \cite{schoenfeld2006survival} \\
        Competing risks & Discharge alive is informative, not censoring &
        Resche-Rigon \cite{rescherigon2006competing};
        Heyard \cite{heyard2019dynamic} \\
        Cox as weak comparator & PH known to fail in sepsis; quantify what the
        field's default costs &
        Kasal et al.\ \cite{kasal2004cox} \\
        No case--control matching & Control alignment inflates apparent
        performance &
        Moor et al.\ \cite{moor2021review} \\
        Treatment-anchored eval & Models learn clinician suspicion, not
        deterioration &
        Kamran et al.\ \cite{kamran2024evaluation} \\
        Anchor-independent label & Sepsis-3, SEP-1 and CDC ASE all key on
        treatment timing, so a within-family comparison cannot separate
        definitional from empirical label variance &
        Kamran \cite{kamran2024evaluation};
        Dutta \cite{dutta2026performance} \\
        Time-based split primary & Sepsis is the worst-affected MIMIC-IV task
        under temporal shift; no algorithmic fix &
        Guo et al.\ \cite{guo2022temporal} \\
        Comparators refitted in-pipeline & Published rankings carry no
        out-of-distribution information &
        Reyna et al.\ \cite{reyna2020utility} \\
        Utility Score primary & AUROC survives the transition that inverts
        clinical value &
        Wang et al.\ \cite{wang2025srpm} \\
        Utility swept, not fixed & A threshold-dependent metric at one cut-off
        measures calibration, not value &
        Reyna \cite{reyna2020utility};
        Wong \cite{wong2026epic} \\
        Calibration co-primary & Threshold-driven instruments are
        calibration-dependent whatever their AUROC &
        Van Calster et al.\ \cite{vancalster2019calibration} \\
        Alert burden reported & 1.4 false alerts per true alert is the
        reference; worse is disqualifying, not a limitation &
        Moor et al.\ \cite{moor2023review} \\
        NEWS2 as benchmark & qSOFA is guideline-deprecated for screening; a
        LASSO-derived score has already failed against NEWS2 &
        Evans \cite{evans2021ssc};
        Mellhammar \cite{mellhammar2020news2} \\
        GBT comparator & The statistics-versus-ML equivalence is contested
        precisely on temporal tasks &
        Christodoulou \cite{christodoulou2019cox};
        Fagerstr{\"o}m \cite{fagerstrom2019lisep} \\
        Subgroup label sensitivity & Label composition varies demographically;
        onset-model subgroup discrimination is unreported &
        Wang et al.\ \cite{wang2022equity} \\
        TRIPOD+AI / PROBAST+AI & The reporting standard covers regression and
        machine learning under one checklist &
        Collins \cite{collins2024tripodai};
        Moons \cite{moons2025probastai} \\
\end{longtable}
}
